\section{Foundations of Agentic Systems}
\label{sec:background}

The term \emph{agent} predates large language models by decades. Classical artificial intelligence defines agency through a control loop: an entity perceives an environment, maintains information relevant to an objective, and selects actions that alter what happens next~\cite{russell2021aima}; multi-agent systems additionally make autonomy, interaction, and coordination first-class concerns~\cite{wooldridge2009multiagent}. The LLM era changes the representational substrate, not this contract. Natural language can encode goals, observations, plans, tool requests, and messages, but those strings become reliable system state only when mapped into explicit, typed, and persistent records. An LLM is therefore best understood as a policy and representation transformer within an agent architecture, not as the architecture itself. A modern agentic system must preserve what was observed, which evidence was accepted, which goal is active, what action was proposed, why it was permitted, and what result followed. That data model is the bridge from evidence to bounded action. This section uses reactive control, deliberative planning, decision-making under uncertainty, hierarchical organization, and multi-agent coordination to explain which component should own each transition and under what constraints.

\subsection{An operational model of agency}

We model a bounded agent as a tuple
\begin{equation}
\mathcal{A}=\langle \mathcal{O},\mathcal{S},\mathcal{G},\mathcal{U},\pi,\mathcal{V},\mathcal{B}\rangle,
\end{equation}
The tuple separates representations from control responsibilities. $\mathcal{O}$ specifies the admissible observation and evidence types; $\mathcal{S}$ is the maintained, persistent state space; $\mathcal{G}$ contains task specifications, including success and termination criteria; and $\mathcal{U}$ defines the schemas for actions or tool calls the architecture can issue. The policy $\pi$ proposes an action from the current state and task, but it does not by itself authorize execution. $\mathcal{B}$ supplies externally enforceable bounds such as permissions, time, cost, source access, and maximum steps, while $\mathcal{V}$ evaluates the transition and determines whether to accept, repair, escalate, or terminate. For $g\in\mathcal{G}$, $\pi$ may be an LLM, a symbolic controller, a learned policy, or a composition of them. This prevents the generative component from defining evidence, state, action, authority, and outcome at once. Observations, goals, proposals, policy decisions, and outcomes must remain distinct typed records so the loop can be inspected or resumed.

\begin{samepage}
At step $t$, an observation updates state, the policy proposes an action, the bounds gate authorizes it, and the environment returns the next observation:
\begin{align}
s_t &= f(s_{t-1},o_t),\\
\tilde{u}_t &\sim \pi(\cdot\mid s_t,g,\mathcal{U}),\\
u_t &= \Gamma_{\mathcal{B}}(\tilde{u}_t,s_t),\\
o_{t+1} &\sim P_{\mathcal{E}}(\cdot\mid s_t,u_t),\\
d_t &= \mathcal{V}(s_t,\tilde{u}_t,u_t,o_{t+1}),
\end{align}
\end{samepage}
The equations separate proposal, authority, and outcome. The update function $f$ first commits the new observation to maintained state. The policy $\pi$ then emits a candidate $\tilde{u}_t$; it does not execute it. The gate $\Gamma_{\mathcal{B}}$ evaluates that proposal against the current state and operational bounds, returning either an authorized $u_t\in\mathcal{U}$ or $\bot$. The latter is a typed result rather than an exception: it records denial, abortion, a controlled no-op, or escalation. Only the authorized action reaches the environment kernel $P_{\mathcal{E}}$, whose stochastic response captures changing APIs, tools, users, and physical or digital environments. Execution success is therefore not equivalent to task success. The validator $\mathcal{V}$ compares the prior state, proposed and authorized actions, and resulting observation to produce $d_t$, such as accept, repair, escalate, or terminate. That decision and its supporting evidence should be persisted with the transition. The formulation imposes two useful invariants: no proposed action bypasses the bounds gate, and no observed outcome bypasses validation. An LLM may implement part of $\pi$ or $\mathcal{V}$, but it cannot supply these guarantees alone; agency emerges from the complete, inspectable loop.

\subsection{Reactive, deliberative, and hybrid control}

Reactive and deliberative control differ principally in how much structured state stands between observation and action. A reactive controller maps the current observation and a compact local state directly to an action. This minimizes latency and can be robust when preconditions are observable and local rules are sufficient, but it provides little machinery for aggregating evidence across time, representing delayed consequences, or comparing alternative futures. A deliberative controller instead persists a goal, an action model, and candidate future states before selecting what to do. STRIPS and PDDL make this representational requirement explicit by encoding actions as transition operators with typed preconditions and effects~\cite{fikes1971strips,mcdermott1998pddl}; BDI architectures separate beliefs, desires, and intentions, distinguishing informational state, objectives, and committed courses of action~\cite{rao1995bdi}. These are not merely historical labels. Reactive control deliberately narrows the representation and decision horizon, whereas deliberation expands both the state that must be maintained and the transitions that must be validated. The appropriate choice depends on the evidence and consequence structure of the task, not on whether an LLM is present.

Most practical LLM agents are hybrid because proposal and execution have different assurance requirements. A language model planner can interpret an underspecified goal, decompose it, and propose or revise a candidate plan; a reactive execution layer must translate each step into a typed command, check its schema, preconditions, permissions, idempotency, and budget, and commit the result to state. Toolformer demonstrates that a model can learn which API to call, when to call it, and which arguments to supply~\cite{schick2023toolformer}; ToolLLM extends tool learning evidence to single- and multi-tool solution paths over a large API collection~\cite{qin2024toolllm}. These results establish proposal competence, not execution authority or postcondition assurance. A plan expressed in language remains a hypothesis about the environment: API state, user permissions, external data, and earlier side effects may invalidate it before the next step. The executor should therefore reread current evidence before every action, bind and validate the concrete parameters, then persist the returned value, error, side effect, and provenance before replanning. Hybrid control allows deliberation to remain flexible while making operational transitions deterministic where possible and auditable where uncertainty remains.

ReAct made the hybrid loop explicit by interleaving reasoning traces, environment actions, and observations~\cite{yao2023react}. Its central architectural contribution is that new environmental evidence can revise the next decision. A textual scratchpad, however, is not a progress record: it can rationalize an error, omit a side effect, repeat an equivalent action, or claim completion without representing the required postconditions. Interactive benchmarks expose this gap. AgentBench identifies long-horizon reasoning, decision-making, and instruction following as recurring obstacles~\cite{liu2024agentbench}, while the original WebArena study reported only 14.41\% end-to-end success for its best GPT-4-based agent, compared with 78.24\% for humans~\cite{zhou2024webarena}. Each loop iteration should therefore update a typed progress ledger containing the intended subgoal, action identifier, checked preconditions, observation, state delta, postcondition verdict, retry count, and remaining budget. Stopping and escalation logic should read that ledger rather than trust the model's assertion that it is finished. A reason--act loop becomes a reliability mechanism only when environment evidence is converted into measurable progress, externally validated transitions, and bounded termination decisions.

\subsection{Decision-making under partial observability}

Partial observability is the normal operating condition of an agent, not an edge case. The system acts on evidence about the environment rather than on the environment's complete state: a user may omit constraints, a web interface may expose only the current view, a tool may return partial or stale data, and another agent may hold private information. Partially observable Markov decision processes (POMDPs) formalize this separation by maintaining a belief over possible states and selecting actions according to both expected task value and the value of additional information~\cite{kaelbling1998pomdp}. An LLM-centered system need not encode that belief as a calibrated probability distribution in every application, but it must preserve an equivalent representational distinction. Observed facts, inferred claims, competing hypotheses, confidence, source, and timestamp should remain separately addressable, and a new observation should revise rather than silently overwrite the prior evidence state. Querying, retrieving, diagnosing, or requesting clarification can then be chosen as information gathering actions. The reliability failure is not that the agent lacks complete state; it is that the architecture lets incomplete evidence masquerade as complete knowledge before action.

The systems consequence is an evidence ledger that keeps observations, inferences, assumptions, and conflicts separately addressable. A scratchpad may assist reasoning, but it cannot be authoritative state because it lacks stable identity, provenance, and update semantics. Information gathering is justified when it reduces decision-critical uncertainty enough to change whether the system proceeds, asks, retrieves, abstains, or escalates. The precise ledger schema belongs to the architecture section; the foundational point is that partial observability becomes governable only when uncertainty persists as state rather than disappearing into fluent prose.

\subsection{Hierarchical control and temporal abstraction}

Hierarchy is a compression of the decision horizon: a controller reasons over subgoals and reusable skills instead of reconsidering every primitive action at every step. Hierarchical task networks formalize this idea by recursively refining an abstract task until its leaves are executable operators~\cite{erol1994htn}; the options framework represents a temporally extended policy through its initiation set, internal policy, and termination condition~\cite{sutton1999options}. Both traditions show that a reusable unit must be represented as more than a natural language label. It requires conditions under which it may start, the state and evidence it consumes, the behavior it controls, and a criterion for recognizing its completion or failure. Modern agentic systems instantiate the same abstraction through planner--executor separation, manager--worker workflows, task graphs, and skill libraries. Voyager, for example, stores increasingly complex behaviors as executable code and improves them using environment feedback, execution errors, and self-verification~\cite{wang2023voyager}. A collection of agents does not automatically create hierarchy: hierarchy exists only when parent and child task states, authority, and result admission are explicitly related. The resulting task tree is a persistent control structure that connects high-level intent to bounded executable actions.

Hierarchy is valuable only when abstraction reduces decision load without erasing control evidence. A delegated result should return as a candidate state transition and pass an admission gate before it changes the parent plan. This boundary can contain local failure, but it also adds routing, synchronization, duplicated context, and recovery costs. MultiAgentBench confirms that topology and planning strategy materially affect milestone completion~\cite{zhu2025multiagentbench}. Hierarchy is therefore a contract between parent and child state (goal, evidence, authority, budget, expected artifact, and completion status), not conversational delegation. Section~\ref{sec:architecture} develops the corresponding task and admission records.

\subsection{Language models as policies over structured actions}

A pretrained LLM produces a probability distribution over text; an executable system requires a controlled choice among typed actions. The tool interface is the boundary that performs this conversion. A model output is interpreted as a proposal naming an operation and supplying arguments, after which a runtime can parse, validate, authorize, execute, and record it. MRKL systems established the routing pattern by selecting external modules for a query~\cite{karpas2022mrkl}; Toolformer showed that a model can learn when to invoke APIs and which arguments to provide~\cite{schick2023toolformer}; ReAct connected action selection to an iterative stream of environmental observations~\cite{yao2023react}. Together, these approaches expand language modeling into policy proposal, but they do not make generated actions correct or safe by construction. Typed function calling narrows the output space and makes malformed calls rejectable; a schema-valid call can still be unsupported by current evidence, inconsistent with state, unauthorized, redundant, or irreversible. The system should therefore preserve both the model's original proposal and its normalized action object so that interpretation errors remain inspectable. Autoregressive generation becomes part of an agent policy only after its output crosses this explicit type boundary and remains subject to runtime control.

The language policy should therefore sit behind a versioned action contract rather than a prose-only tool description. The model proposes an operation and arguments; the runtime validates their schema, evidence, authority, budget, and expected effects before execution and records a typed receipt afterward. BFCL shows that correct function selection and argument construction are distinct capabilities that require executable evaluation~\cite{patil2025bfcl}. Yet even a schema-valid call may be unsafe or unsupported. Dependability is joint: model competence matters, but the interface and enforcing runtime determine which errors can reach the world.

\subsection{Memory, external cognition, and persistent state}

Agent memory is an architecture of typed state classes, not one store or the current prompt. Working state, episodes, semantic evidence, profiles, and procedural skills differ in authority, retention, retrieval, and validation~\cite{zhang2025memorysurvey}. MemGPT illustrates the systems value of moving information between limited context and external tiers~\cite{packer2024memgpt}; AgeMem makes add, update, delete, retrieve, summarize, and filter explicit policy actions~\cite{yu2026agemem}. These mechanisms also expose the central risk: a generated note or retrieved episode can influence later decisions without becoming verified evidence. Memory writes and revisions therefore require provenance, scope, versioning, and an admission decision, while retrieval must preserve authorization and freshness. The architecture section develops these lifecycle operations; here memory matters because it extends the agent's control loop across time and turns persistence itself into a bounded action surface.

\subsection{Multi-agent foundations}

A multi-agent system externalizes decision-making across entities with different observations, objectives, capabilities, or authority~\cite{wooldridge2009multiagent}. This distribution is useful when no single controller should see every datum or hold every permission, when subtasks require genuinely different tools or expertise, or when independent exploration and critique produce evidence that can be compared. Contemporary LLM systems encode that structure in different ways: AutoGen treats coordination as configurable conversations among agents~\cite{wu2023autogen}; CAMEL studies role-conditioned cooperation through agent dialogue~\cite{li2023camel}; and MetaGPT organizes software work through specialized roles and workflow artifacts~\cite{hong2024metagpt}. These designs differ in topology and control, but each makes communication part of the architecture rather than merely additional prose. An agent should have an explicit identity, capability and permission profile, observation scope, memory view, active commitments, and allowed message types. A message should similarly identify sender, recipient, task and correlation identifiers, claim or requested action, evidence references, timestamp, and confidence or commitment status. Receiving a message does not make its content true; an admission rule determines whether it updates shared task state. In this sense, a multi-agent system is a distributed state machine whose transitions happen through typed messages. Splitting one model prompt into several personas without distinct evidence, authority, or state boundaries creates dialogue, not substantive distribution.

The distribution introduced above is useful only when it creates complementary information, capability, or authority. Cloning one model behind several role labels can preserve shared blind spots: agents may repeat the same unsupported premise, while agreement amplifies confidence without adding independent evidence. The empirical record is correspondingly mixed. Iterative multi-agent debate improved factuality and reasoning on the controlled tasks studied by Du et al.~\cite{du2024multiagentdebate}; by contrast, Smit et al. found that current debate strategies did not reliably outperform self-consistency or ensembles of independent reasoning paths without careful tuning~\cite{smit2024mad}. MultiAgentBench further shows that milestone completion changes with coordination topology and planning strategy, while exposing repetition and inefficient collaboration as observable failure modes~\cite{zhu2025multiagentbench}. The design question is therefore not how many agents participate, but what independence the architecture actually creates. A coordination contract should specify the communication topology, admissible message types, ownership of shared state, conflict resolution and aggregation rules, termination criteria, and escalation paths. Roles should be functionally diverse, using different evidence sources, tools, permissions, or verification duties, and their outputs should be merged by an evidence-aware rule that preserves provenance and dissent. Collective judgment must therefore rest on independently grounded contributions and explicit reconciliation, not conversational agreement.

\subsection{A bounded interpretation of autonomy}

Autonomy is not a property of model output; it is an allocation of decision rights across transitions. Human--automation research treats automation as stage-specific~\cite{parasuraman2000levels}, and the NIST AI Risk Management Framework requires explicit oversight roles~\cite{tabassi2023airmf}. An LLM-centered system should therefore represent autonomy as an enforceable profile rather than a label. Capability and authority remain distinct: knowing how to initiate an action does not grant permission to execute it. Relevant dimensions include:
\begin{itemize}
    \item \textbf{decision autonomy}: which branches or actions the model may choose;
    \item \textbf{execution autonomy}: whether proposed actions are executed automatically;
    \item \textbf{temporal autonomy}: how long the system may operate before review;
    \item \textbf{resource autonomy}: what compute, money, tools, and data it may consume;
    \item \textbf{social autonomy}: whether it may communicate or delegate to people and agents;
    \item \textbf{adaptation autonomy}: whether it may alter memory, prompts, policies, or code.
\end{itemize}
These dimensions form a profile, not a ladder. A research agent may search and synthesize independently while remaining unable to send messages, modify records, or spend money without approval. High automation can coexist with high human control~\cite{shneiderman2020hcai}; what matters is the separation among capability, delegated authority, and observed behavior. Governance and evaluation should compare those quantities directly, because labels such as ``copilot'' and ``fully autonomous'' conceal the transitions that actually require control.

\subsection{Conceptual analogies and their limits}

Neuroscientific and cognitive theories are useful as sources of architectural hypotheses, not as evidence that an LLM-centered system implements the inspiring mechanism. Marr's computational, algorithmic, and implementation levels provide a useful test~\cite{marr1982vision}: a metaphor needs a precise mapping, an architectural claim needs an executable mechanism and ablation, and a biological claim needs causal evidence. The Thousand Brains Theory may motivate specialized evidence views and reconciliation~\cite{hawkins2021thousandbrains}, but specialized modules are not thereby cortical columns. We retain only the operational lesson that multiple evidence-bearing representations can be reconciled before bounded action.

The foundations converge on one control structure: partial evidence becomes typed state; state supports a proposal; goals, permissions, and budgets constrain authorization; and observed effects update the next decision. Memory carries the loop across time, hierarchy across temporal scales, and multi-agent messaging across responsibility boundaries. The common requirement is a representation contract for observations, claims, plans, authority, effects, and failure. The next section asks where those records live, who owns them, and which component enforces their transitions.
